Figure~\ref{fig:emw_levels_rx_combined} shows level outcomes in the stacked treatment and control groups. Before coverage, obesity-indicated GLP-1 prescriptions were near zero in both treated and control states. Coverage increased GLP-1 obesity prescribing immediately and the size of the effect grew over time. Figure~\ref{fig:emw_rx_combined} presents event study estimates of the effect of Medicaid GLP-1 obesity coverage on prescriptions per 100 enrollee-months. In the 10-state design, obesity GLP-1 prescriptions per 100 enrollee-months increased by \rxObTenZeroEmw{} (SE = \rxObTenZeroEmwSE{}) in the quarter of adoption, and grew to \rxObTenSixEmw{} (SE = \rxObTenSixEmwSE{}) six quarters later. The 7-state design, which allows longer follow-up, shows continued growth: prescriptions per 100 enrollee-months increased by \rxObSevenZeroEmw{} (SE = \rxObSevenZeroEmwSE{}) in the quarter of adoption, and reached \rxObSevenNineEmw{} (SE = \rxObSevenNineEmwSE{}) nine quarters after adoption. Coverage for GLP-1 obesity drugs had no effect on diabetes or cardiovascular-indicated GLP-1 prescriptions in either design; point estimates are close to zero and statistically insignificant at every post-treatment event time. This finding is consistent with physicians not having previously prescribed GLP-1 obesity medications off-label under diabetes indications to circumvent coverage exclusions: if such off-label prescribing had been common, we would expect diabetes-indicated prescriptions to decline as obesity coverage provided a direct channel.

Gross Medicaid reimbursements for obesity-indicated GLP-1 drugs closely track prescriptions. Figure~\ref{fig:emw_levels_spending_combined} shows quarterly gross spending per 100 enrollee-months in the treated and control states, and Figure~\ref{fig:emw_spending_combined} presents the corresponding event studies. These spending figures reflect Medicaid reimbursements before manufacturer rebates. In the 10-state design, coverage increased gross reimbursements per 100 enrollee-months by \$\spObTenSixEmw{} (SE = \$\spObTenSixEmwSE{}) six quarters after adoption. In the 7-state design, coverage increased gross reimbursements per 100 enrollee-months by \$\spObSevenFourEmw{} (SE = \$\spObSevenFourEmwSE{}) after four quarters, and by \$\spObSevenNineEmw{} (SE = \$\spObSevenNineEmwSE{}) after nine quarters. There were no significant effects on spending for diabetes or cardiovascular-indicated GLP-1 drugs. Total gross Medicaid reimbursements per 100 enrollee-months for all GLP-1 drugs (obesity, diabetes, cardiac) increased by \$\spAllTenSixEmw{} (SE = \$\spAllTenSixEmwSE{}) at six quarters in the 10-state design and by \$\spAllSevenNineEmw{} (SE = \$\spAllSevenNineEmwSE{}) at nine quarters in the 7-state design, driven almost entirely by the obesity indication.

Under the Medicaid Drug Rebate Program (MDRP), manufacturers pay rebates to state Medicaid programs, so net spending is lower than the gross reimbursements reported in the SDUD. We estimate net spending by applying statutory MDRP rebate floors to drug-level price data (see Appendix~\ref{sec:appendix_b} for details). Figure~\ref{fig:emw_levels_net_spending_combined} shows net spending levels and Figure~\ref{fig:emw_net_spending_combined} presents net spending event studies. In the 10-state design, the effect on net obesity spending per 100 enrollee-months was \$\netSpObTenSixEmw{} (SE = \$\netSpObTenSixEmwSE{}) six quarters after adoption, compared to \$\spObTenSixEmw{} in gross reimbursements---a reduction of approximately \spReductionObTenSixEmw{} percent. In the 7-state design, net obesity spending per 100 enrollee-months was \$\netSpObSevenNineEmw{} (SE = \$\netSpObSevenNineEmwSE{}) nine quarters after adoption, compared to \$\spObSevenNineEmw{} in gross reimbursements---a reduction of approximately \spReductionObSevenNineEmw{} percent. MDRP rebates thus reduce the estimated spending effect by roughly one quarter. The pattern for diabetes or cardiovascular-indicated GLP-1 net spending mirrors the gross results: no significant effects in either design.

The pre-treatment event study coefficients support the parallel trends assumption. Pre-treatment coefficients are quantitatively small for all outcomes. For obesity-indicated GLP-1 outcomes, pre-treatment coefficients are close to zero in substantive terms but statistically significant, because pre-treatment prescriptions are near zero in both treated and control states, residual variance is very small, and even trivial differences reach significance.

Although not substantively important, the pre-treatment patterns show that obesity-indicated GLP-1 utilization is very low but not literally zero before coverage: 9 of 10 treated states and 30 of 33 control states have some pre-treatment prescriptions, mostly filled through MCOs. This suggests that some MCOs may offer restricted coverage even when the state formulary does not.\footnote{California is an exception: it had some pre-treatment utilization from both FFS and MCO claims.}

To gauge the aggregate fiscal implications of the per-enrollee estimates, we scale the ATTs from the 7-state enrollee-month weighted design by each state's September 2025 Medicaid enrollment to project monthly users and annual net spending. We report two benchmarks: an ``average post-period'' estimate based on the mean ATT across event times 0--9, and an ``event time 9'' estimate reflecting more mature utilization nine quarters after adoption. Table~\ref{tab:budget_impact_covering} presents projections for the \numCoveringStates{} states that adopted coverage during our study period. At the average post-period benchmark, these states would have approximately 104,600 enrollees filling a GLP-1 obesity prescription each month, generating \$1.18 billion in annual net Medicaid spending.\footnote{The enrollee-month weighted design is our main specification; see Section~\ref{sec:methods} for details.} At event time 9, projected utilization grows to roughly 243,600 monthly users and \$\netCostCoveringEtNine{} in annual net spending. California alone accounts for \$\netCostCAetNine{} of the event-time-9 total, driven by its \enrollCA{} Medicaid enrollees.

Table~\ref{tab:budget_impact_noncovering} extends the projections to the 33 states that had not adopted coverage by mid-2025. If these states were to cover GLP-1 obesity drugs and experience similar per-enrollee effects, projected annual net spending would be \$1.61 billion at the average post-period benchmark and \$\netCostNoncoveringEtNine{} at event time 9. The largest non-covering states drive much of the total: New York (\$530 million at event time 9), Texas (\$347 million), and Florida (\$314 million). Combining all 50 states, projected annual net Medicaid spending on GLP-1 obesity drugs ranges from approximately \$2.8 billion at the average post-period benchmark to roughly \$6.3 billion at event time 9. These projections assume uniform per-enrollee effects across states and do not account for differences in demographics, formulary design, or supplemental rebates.

The SDUD results show that Medicaid coverage increased GLP-1 utilization and spending within the program, but a separate question is whether this coverage displaced private out-of-pocket spending on GLP-1 products. Figure~\ref{fig:ce_national_timeseries} uses Consumer Edge transaction data to document the rapid growth of a parallel private market. National spending on online GLP-1 compounding platforms---telehealth companies selling compounded semaglutide and tirzepatide---grew roughly 35-fold between 2019 and 2026, while Lilly Direct, a manufacturer direct-to-consumer channel launched in late 2024, reached spending levels comparable to those of the compounding platforms within months. In contrast, spending on weight loss platforms such as Noom and WeightWatchers declined over the same period. If Medicaid coverage of manufacturer-produced GLP-1 drugs crowded out private spending---for example, because some Medicaid-eligible individuals were previously paying out of pocket for compounded GLP-1 products---we would expect online GLP-1 platform transactions to fall in states that adopt coverage. Figure~\ref{fig:emw_ce_compounder_combined} shows online GLP-1 spending levels in the treated and control groups and presents event study estimates. Post-treatment point estimates are negative at every event time in both designs, directionally consistent with crowd-out. In the 10-state design, the estimated effect on online GLP-1 spending per 100 enrollee-months was $-$\$\ceCompTenSixEmw{} (SE = \$\ceCompTenSixEmwSE{}) six quarters after adoption. In the 7-state design, the effect was $-$\$\ceCompSevenNineEmw{} (SE = \$\ceCompSevenNineEmwSE{}) nine quarters after adoption. None of the estimates are statistically significant, and pre-treatment coefficients show no differential trends (joint F-test $p > \cePrePMin{}$ for all outcome-design combinations). Lilly Direct launched too recently for event study analysis in either design. The negative point estimates are directionally consistent with some displacement of private spending, but the effects are imprecisely estimated and do not support a strong crowd-out conclusion.\footnote{Even taking the point estimates at face value, the implied displacement would be less than 2 percent of net Medicaid spending on obesity-indicated GLP-1s, indicating that coverage generated predominantly new utilization.}

